\chapter{Results} % (fold)
\label{chap:results}

The role of chirps as communication signals of wave-type weakly electric fish
is still mysterious, partly because they are not easy to study. I have now
introduced two neural networks that can be used to detect and assign chirps to
their emitter. In the following chapter, we see how these models can be
evaluated and integrated into a pipeline to automate chirp detection. This will
allow us to detect chirps on a large dataset from a behavioral experiment to
explore relationships between chirps, dominance status, and agonistic
interactions in competitive contexts of \textit{A.~leptorhynchus}.

\section{Architecture of the chirp detection pipeline} % (fold)
\label{sec:the_pipeline}

The chirp detection pipeline includes two feature extraction and two inference
phases. Initially, spectrograms are extracted from EOD recordings and are
passed into the detection model that returns time and frequency ranges, where
it detected chirps. Following chirp detection, a second feature extraction
phase gathers the temporal evolution of the powers of potential emitter
EOD$f$s, which a subsequent model classifies in the final inference step. This
way, each detected chirp is assigned to a fish. The full pipeline can be
accessed as a \texttt{Python} package and all code necessary to detect chirps
is freely available on GitHub
(\url{https://github.com/weygoldt/chirpdetector}). The upcoming sections detail
all processing steps and the models involved.

\subsection{Detection feature extraction}

The recordings from electrode grids have as many channels as there are
electrodes. The initial feature extraction phase transforms raw electric
recordings into spectrograms using a sliding window approach on every channel.
These spectrograms are then combined into a sum spectrogram and converted to
the decibel scale. Optimized for handling large datasets, these spectrograms
are computed \enquote{on-the-fly} and discarded post-detection model analysis.
To facilitate detection of small objects by the model and manage computational
costs effectively, spectrograms are segmented into smaller tiles
\parencite{huangTilingStitchingSegmentation2019}. Each tile spans
\SI{15}{\second} with \SI{1}{\second} overlap, and covers a frequency range of
\SI{1}{\kilo\hertz} with a \SI{500}{\hertz} overlap up to a maximum frequency
of \SI{2}{\kilo\hertz}. For instance, a minute-long recording yields four
time-axis windows and three frequency-dimension tiles, totaling 12 spectrograms
processed concurrently.

These spectrograms are then analyzed by the detection model to identify chirp
bounding boxes, with parallel processing capacity constrained only by GPU
memory. On a NVIDIA GeForce RTX 4080 with 16 GB of graphical memory, the system
can process three minutes of data from 16 electrodes sampled at
\SI{20}{\kilo\hertz} simultaneously. All performance-relevant parameters are
adjustable depending on the available hardware.

\subsection{Chirp detection}

Inference is performed using a YOLOv8 model, encapsulated by a wrapper to
standardize output formats, such as bounding boxes and confidence scores,
across different model architectures. This wrapper, based on an abstract base
class, facilitates easy model interchangeability and allows for the
incorporation of new models as they become available.

I tested the performance of all YOLOv8 versions as well as YOLOv9E. Both
training and cross-validation indicated that YOLOv8XL outperforms others in
precision, recall, and F1 score (figure \ref{fig:detector_performance} A-C).
Precision, informative about how many of the  predicted chirps are actually
real chirps, (equation \ref{eq:ap}), reaches above 99\% for the YOLOv8XL
(figure \ref{fig:detector_performance} A). Recall, relating to how many of the
chirps that are truly present in the dataset, are also detected (equation
\ref{eq:ap}), also surpasses 99\% for YOLOv8XL (figure
\ref{fig:detector_performance} B). Together, this shows that most of our
predictions are trustworthy and that most chirps in the dataset will be
detected. Only approximately 1\% of detections will be false positives and only
approximately 1\% of existing chirps will be missed. Optimal detection
thresholds to balance precision and recall are identified at around 50\% for
all models by the F1 score (figure \ref{fig:detector_performance} C). YOLOv8XL
was leading in average precision that reached up to 86\% (figure
\ref{fig:detector_performance} D, E, equation \ref{eq:ap}). The average
precision is a measure combining precision and recall into a single metric to
make comparisons of models easier. This is lower to the previous metrics as
here, the average precision is evaluated at AP@[.5:.01:.95] while the curves
are evaluated at a single IoU threshold of 0.5 (see section \ref{sec:deteval}).
Still, AP@[.5:.01:.95] of over 85\% can be regarded as a very precise object
detector, surpassing the current state of the art benchmarks for multi-class
detectors \parencite{wangYOLOv9LearningWhat2024}. Surprisingly, YOLOv9E, shows
subpar performance despite beeing the newest model, potentially due to training
instabilities (figure \ref{fig:yolo_crossval_training}). The performance
metrics of YOLOv8M to XL are closely matched, with average precisions exceeding
85\%. These medium sized models may suffice for less powerful devices or
embedded systems.

\begin{figure}[ht!]
  \begin{center}
    \includegraphics[width=1\textwidth]{../figures/kfold_modelsweep_performances.pdf}
  \end{center}
  \caption{
    Performance metrics for detection models on the validation dataset,
    including precision, recall, and F1 score. The F1 score serves as a
    performance indicator to determine the optimal confidence thresholds.
    Results, shown as median values from 5-fold cross-validation with
    \ordinalnum{25} and \ordinalnum{75} percentiles, though variations are
    minimal. Plots focus on areas where model performances diverge. Model
    precision quickly stabilizes near 100\%, with YOLOv8N slower to plateau
    (A). Model recall remains steady up to a 0.8 confidence, then declines (B).
    YOLOv8M-XL exhibit higher recall, and are less affected by the confidence,
    with YOLOv8XL leading. Optimal F1 scores occur around a confidence level of
    0.5 for all models (C). Larger models score higher, with YOLOv8XL achieving
    the top F1 score. Precision-recall curves, based on an IoU threshold of
    0.5, largely overlap, with YOLOv8XL showing the greatest area under the
    curve (D). Average precision, calculated across IoU thresholds from 0.5 to
    0.95, indicates overall model efficacy (E). Larger models outperform, with
    YOLOv8XL leading in average precision. YOLOv9E performs comparatively
    poorly and exhibits strong variability, probably due to the instabilities
    during training illustrated in figure \ref{fig:yolo_crossval_training}.
  }
  \label{fig:detector_performance}
  \vspace{0.5cm}
\end{figure}

Detected bounding boxes undergo several post-processing steps: Initially,
low-confidence detections are filtered out through thresholding. Subsequently,
non-maximum suppression eliminates overlapping boxes using the \acrshort{IoU}
metric, effectively reducing false positives and ensuring each chirp is
represented by a single bounding box. These bounding boxes mark detected
chirps, but we still do not know, which of the fish emitted these chirps. This
is why we need a second processing stage that assigns the detected chirps to
the emitting fish.

\subsection{Feature extraction for chirp assignment}

The second feature extraction, applied to each detected bounding box, begins by
estimating the fundamental frequency of all the fish potentially emitting the
chirp. This involves first interpolating the spectrogram within the bounding
box to $100 \times 100$ samples and calculating the power spectral density by
summing the resulting matrix along the time axis. Peaks in this spectrum
indicate the fundamental frequencies of potential emitters. A \SI{10}{\hertz}
window is used around each peak to sum power values in the spectrogram across
the frequency axis. The result is array representing the power of the
fundamental frequency over time for each potential emitter within a bounding
box. This process identifies deviations in the power of the fundamental
frequency associated with chirp emissions, where the power at the baseline
frequency drops. In non-chirping instances the power remains approximately
constant. The whole process is visualized in figure
\ref{fig:assignment_feature_extraction}. This feature extraction is efficient,
leveraging already computed spectrograms rather than having to go back to raw
electric recordings.

\begin{figure}[t]
  \begin{center}
    \includegraphics[width=\textwidth]{../figures/assignment_training.pdf}
  \end{center}
  \caption{
    Training curves and performance evaluation of the assignment model across
    5-fold cross-validation. The lines correspond to the median across the
    cross-validation folds and the colored areas to the \ordinalnum{25} and
    \ordinalnum{75} percentile respectively. The training as well as the
    validation loss converges after 11 epochs (A). F1-scores peak at a model
    confidence of approximately 82\% (B). The average precision of the model,
    i.e., the area under the precision-recall curve reaches approximately 97\%
    (C). The AUC of the ROC curve (D) reaches 99\% but a ROC evaluation has
    been shown to be overly optimistic when datasets are imbalanced
    \parencite{brancoSurveyPredictiveModelling2015}.
  }
  \label{fig:assignment_training}
  \vspace{0.5cm}
\end{figure}

\subsection{Chirp assignment}

For each bounding box, features are passed to the assignment model, which
evaluates the probability of each power array representing a chirp emitted by a
fish. The array with the highest probability identifies the emitter, assigning
its frequency as the fundamental frequency at that moment. The assignment
process, encapsulated by a wrapper, facilitates feature extraction and model
inference, with an abstract base class defining the model interface. Similar to
the detection model, this setup allows for easy model interchangeability and
simplifies the integration of new models into the pipeline.

The assignment model demonstrates high accuracy, achieving an F1-score of 91\%,
average precision of 97\%, and a AUC of 99\%, indicating effective chirp
assignment in most instances (figure \ref{fig:assignment_training} B-D). To
better understand the model's performance, I visualized the decision boundary
by reducing the 100-dimensional power traces to two dimensions using a simple
auto-encoder model (figure \ref{fig:assignment_zspace}). This reduction was
compared to results from \acrshort{PCA}, showing similar outcomes. The
visualization reveals a transitional zone between arrays classified as chirp
emitters and non-emitters, where ambiguous cases, likely representing short
chirps with less noticeable power drops, are found. Chirps with distinct power
reductions are clearly separated. The decision boundary, depicted with colors
for intermediate prediction scores, resembles a parabola rotated to
\SI{45}{\degree}. Most cases fall exactly onto the decision boundaries as
illustrated by the marginal distributions. This is surprising, given the high
performance metrics of the model. It indicates that the two-dimensional
representation is likely too simple to represent the true multidimensional
decision boundary and supports the used of a neural network, as a conventional
machine learning classifier running merely on the principle components probably
would have failed. But this remains to be verified.

\begin{figure}[!t]
  \begin{center}
    \includegraphics[width=\textwidth]{../figures/z_space_assignment_nn.pdf}
  \end{center}
  \caption{
    A visual simplification of the CNN classifier's decision boundary to a
    two-dimensional view for the evaluation dataset. Each point represents a 2D
    latent space mapping of power over time for a potential chirp emitter
    frequency, derived using a fully connected auto-encoder neural network
    trained on 20\% of the dataset. Points are colored based on the binary
    classifier's score, revealing a clear emitter class cluster that
    transitions into the non-emitter class. Insets highlight specific samples
    and their 2D mappings, illustrating a clear separation with definite
    emitters (top right) and non-emitters (bottom left) based on peak presence.
    Marginal distributions on each axis correspond to the chirp emitters and
    non-emitters after thresholding, based on an ideal 82\% score threshold
    from the F1-curve (figure \ref{fig:assignment_training}). These show
    overlapping distributions at the decision boundary. Despite this overlap,
    the classifier achieves an impressive 97\% average precision, as detailed
    in figure \ref{fig:assignment_training}.
  }
  \label{fig:assignment_zspace}
  \vspace{0.5cm}
\end{figure}

\newpage
\subsection{Postprocessing}

The two-step inference pipeline returns bounding box coordinates and predicted
fundamental frequencies of chirp-emitting fish. These outputs are integrated
with the \texttt{wavetracker} package, which tracks fish fundamental
frequencies over time on spectrograms and assigns identities to each frequency
track \parencite{raabAdvancesNoninvasiveTracking2022}. This allows
for estimating the number of recorded fish, approximating their positions, and
maintaining their identities. By matching the detected chirps' fundamental
frequencies with the closest fish frequency at the chirp time, chirps are
attributed to specific fish. Detections not aligning with a frequency track,
such as chirps detected on harmonics, are discarded.

The refined dataset contains chirp times, their dimensions on the spectrogram,
and the emitting fish's identity. Analysis of this dataset from staged
competitions revealed notable temporal patterns between chirps and aggressive
interactions, discussed in the following section.

\section{Evaluation of the detected chirps}

The chirp detector was applied to the Raab dataset, as published in
\textcite{raabElectrocommunicationSignalsIndicate2021}, along with additional
post-publication competition trials. This encompassed 69 trials in total, where
two fish competed for shelter over six hours in each trial, including three
hours of their active phase. To enhance signal clarity, a conservative
detection threshold of 0.7 was chosen, likely increasing the number of missed
chirps but minimizing false positives (figure \ref{fig:detector_performance}).
Across a total of 17.25 days of competition data, the detector identified
64,050 chirp instances.

Before exploring the potential role of chirps in the competitions documented by
\textcite{raabElectrocommunicationSignalsIndicate2021}, let us first assess
chirp detection effectiveness in new data through some examples. The detection
system successfully identifies and assigns most chirps, even when they
overlap other baseline EOD$f$ frequencies (figures \ref{fig:chirp_detections_simultaneous}
and \ref{fig:chirp_detections_bigchirps}). However, challenges arise in
specific scenarios. For instance, figure
\ref{fig:chirp_detections_simultaneous} illustrates a scenario where rapid
chirp production by one fish leads to detection failures and unreliable
assignments due to the spectrogram's limited temporal resolution. While such
high chirp rates are uncommon, they could hold significant behavioral interest.
Additionally, detecting chirps produced simultaneously by two fish poses
difficulties when their frequency difference is small, which causes chirps to
merge on the spectrogram. This often results in neural networks misidentifying
them as a single, large chirp. Despite the initial assumption that those cases
must be rare, figure \ref{fig:chirp_detections_simultaneous} presents two
instances of simultaneous chirps; one correctly assigned, and the other
falsely detected as a single chirp.

\begin{figure}[t]
  \begin{center}
    \includegraphics[width=\textwidth]{../figures/chirp_detections_simultaneous.pdf}
  \end{center}
  \caption{
    Detected and assigned chirps on an example spectrogram. Black lines
    correspond to the tracked fundamental frequencies of two fish. White boxes
    surround all detected chirps. The predicted fundamental frequency of the
    chirp emitter is illustrated by the dot in each bounding box. Overall, most
    chirps are detected well. The pipeline fails when chirps happen in too
    rapid successions to be separable on the spectrogram or when fish chirp
    simultaneously. The spectrogram was estimated using the same parameters
    used for chirp detection (nfft=4096, hop length=410) and interpolated
    (bicubic).
  }
  \label{fig:chirp_detections_simultaneous}
  \vspace{0.5cm}
  % 
\end{figure}

\begin{figure}[!h]
  \begin{center}
    \includegraphics[width=\textwidth]{../figures/chirp_detections_bigchirps.pdf}
  \end{center}
  \caption{
    Assignment can completely fail for large and long chirps and produces false
    assignment for large and short chirps. In one case, a large chirp is not
    detected at all. But the majority is detected and assigned correctly. The
    spectrogram was computed as described in figure
    \ref{fig:chirp_detections_simultaneous}.
  }
  \label{fig:chirp_detections_bigchirps}
  \vspace{0.5cm}
  % 
\end{figure}

Detecting large chirps presents another challenge. Although the training
dataset emphasizes these through sample duplication in histogram equalization,
long chirps are still accurately detected (figure
\ref{fig:chirp_detections_bigchirps}). However, assigning these chirps to the
correct sender often fails. This issue may arise because long chirps gradually
return to the baseline EOD$f$, which falls outside the detected bounding boxes.
Since emitter frequencies are determined from the power spectrum within these
boxes, long chirps might not exhibit distinct peaks at the emitter's EOD$f$.
Essentially, the power dip at the baseline is so broad that the baseline
frequency component becomes too faint to identify.

Another issue remains assignment in general: While most large chirps are
correctly assigned despite strongly overlapping with the baseline EOD$f$ of
another fish in figure \ref{fig:chirp_detections_bigchirps}, it includes three
examples of false assignments. Reasons could be amplitude troughs that
coincidentally occur due to noise or movements of the bottom fish, and that
resemble the troughs that happen during chirps. To conclude, while the great
majority of chirps are detected well, small issues remain. In most cases
however, they appear to be fixable.

Now let us take a look at all the bounding boxes that were predicted on the
dataset and are visualized in figure \ref{fig:all_boxes}. Predominantly, these
bounding boxes encapsulate small chirps, with the majority measuring below
\SI{200}{\hertz} in height and spanning \SI{100}{\milli\second} to
\SI{200}{\milli\second} in width (figure \ref{fig:all_boxes} B1). Remember that
the true chirp height and width does not directly map well to the dimensions on
the spectrogram: The most commonly observed small chirp is usually just
\SI{20}{\milli\second} in width but will be much wider on the spectrogram,
dependent on the parameters the spectrogram was computed with. This is the
reason of the compression of the width axis and the lack of chirps below
\SI{100}{\milli\second} in width.

Consistent with the training data in figure
\ref{fig:chirp_overview_clustered_plot}, the vast majority of chirps fall below
\SI{200}{\hertz} in height, as shown by the marginal distribution in figure
\ref{fig:all_boxes}. Another pattern that we have already observed in the
training data is that larger chirps tend to have stronger post-chirp
undershoots (figure \ref{fig:all_boxes} B2). But this pattern only extends to
short chirps, as the few long chirps that were detected do not show such strong
troughs (figure \ref{fig:all_boxes} B4).

In addition, two areas are visually separated from the main point cloud: Area
B3 is shaped similarly to area B2 but contains much less data points and they
appear to be moved up by \SI{100}{\hertz}. The simplest explanation would be
that these are large chirps produced by the fish with the higher EOD$f$, which
got falsely assigned because they overlapped with the frequency of the fish
with the lower EOD$f$. This way, the chirp height is incremented by the
frequency difference between the two fish, which often lies in a range between
\SI{50}{\hertz} and \SI{100}{\hertz}. Three examples for such false assignments
of large chirps are indicated in figure \ref{fig:chirp_detections_bigchirps}.
The second indicator supporting this assumption is the fact that this point
cloud has weaker post-chirp undershoots compared to area B2 in figure
\ref{fig:all_boxes}. This would also be explainable by the false assignment:
The trough is computed by subtracting the predicted baseline by the lower edge
of the bounding box. If a fish with a lower EOD$f$ is falsely assumed to be the
emitter, this difference becomes smaller.

\begin{figure}[!t]
  \begin{center}
    \includegraphics[width=\textwidth]{../figures/plot_chirp_bbox_distributions.pdf}
  \end{center}
  \caption{
    Statistics of all detected bounding boxes on the Raab dataset. Bounding
    boxes are centered at the baseline EOD$f$ of the fish and colored by
    density of the box type (A). Densities are estimated by Gaussian kernel
    density estimate on the first principle component of the height, width and
    trough parameters. The great majority of chirps are surrounded by bounding
    boxes with a width of \SI{100}{\milli\second} to \SI{200}{\milli\second}
    and a height below \SI{200}{\hertz}, which corresponds to small or type 2 
    chirps (B1). Post-chirp undershoots of the frequency become increasingly
    pronounced particularly for large but narrow chirps (B2). Chirps with a
    width beyond \SI{250}{\milli\second} are sparsely represented (B4). Beyond
    a height of \SI{500}{\hertz}, points appear to be a part of the point cloud
    below but incremented along the height dimension by approximately
    \SI{100}{\hertz}, which could be due to erroneous EOD$f$ predictions (B3).
    Chirps at around \SI{200}{\hertz} in height include a surprisingly narrow
    horizontal band that extends to durations of up to almost
    \SI{400}{\milli\second} (B5). These are most likely bounding boxes around
    chirps that happened in rapid succession, so that reliable chirp separation
    on the here used spectrogram parameters was not possible anymore. Chirps
    at a height range corresponding to this band are the only ones observed to
    occur in such rapid succession in this dataset, which would support this
    assumption.
  }
  \label{fig:all_boxes}
  \vspace{0.5cm}
\end{figure}

The second visually distinct area corresponds bounding boxes that are at
approximately \SI{200}{\hertz} in height but form a distinct band that extends
along the width axis (figure \ref{fig:all_boxes} B5). These are most likely
chirps that happen in such a rapid succession that they are inseparable on the
spectrogram. Given the width of \SI{100}{\milli\second} of most chirps in
figure \ref{fig:all_boxes}, rapid chirps that have a smaller inter-chirp
interval this, cannot be separated on these spectrograms. The frequency range
in which they occur supports this assumption, as chirps that are larger than
\SI{200}{\hertz} have not been observed to happen in such rapid succession
here. In some cases, the network still detects them as chirps but detection
boundaries can become much wider as they usually encompass multiple chirps at
once. Another interesting feature of this area is that it contains many chirps
with strong undershoots. This is especially true for more realistic (smaller)
chirp widths in this range. The most likely explanation is, that these are big
chirps produced by fish with a lower EOD$f$, which got falsely assigned to the
fish with the higher EOD$f$. This would drastically increase the trough while
decreasing the chirp height. Usually, the post-chirp trough is not stronger
than $1/3$ of the chirp height. In this area, the troughs are larger than the
chirp heights, which is not realistic. A simple filter that limits the
\enquote{chirp height to trough ratio} could remove these approximately 8,000
suspicious data points. Despite these two areas, most of the 73,134 points fall
into areas that are expected. In combination with the performance metrics, the
two example spectrograms, and hundreds of additional hand-checked spectrograms
the pipeline seems to detect chirps well enough to use the detections to
conduct behavioral analyses.

Now that we have demonstrated that the pipeline not only produces promising
performance metrics but that these metrics are also reflected on real,
previously unseen data, we can move on to analyze which fish produced chirps
and when they produced them.

% \newpage
\section{Subordinates chirp the most}

\textcite{raabElectrocommunicationSignalsIndicate2021} staged competitions
between two fish for access to a shelter. During these competitions, fish
partook in ritualized fights and chasing. In addition to video recordings, they
also recorded the electric fields of both fish using electrode grids. As video
recordings are not available for a large part of the dataset, and not annotated
yet for another part of the dataset, only 19 unique pairings between two fish
have annotated competitive behaviors. Figure
\ref{fig:competition_trial_overview} shows an overview of the available data
across one of these competition trials. A single trial consisted of three hours
of darkness, where fish were active and usually competed. During the subsequent
three illuminated hours, fish usually became inactive and sought shelter. The
fish that stayed in the superior shelter at the onset of the inactive period
was decided to be the winner of the competition or the \enquote{dominant}. In
the example trial, particularly the competition loser produced chirps, all of
which appear to visibly correlate with the chasing events.

To quantify whether this observation is systematically observable, I compared
the total number of chirps produced by dominants and subordinates. Overall,
winners chirped much less compared to competition losers (p<0.0001, d=66.5), a pattern that could be observed for all trials (figure \ref{fig:stats} A).
Losers produced chirps in an average of 29\% $\pm$ 23\% of all chasing
events.

As suggested in the literature, chirps could have been used as a submissive
signal by competition losers to stop extended chasing by the competition winner
or decrease the probability of future attacks by winners
\parencite{henningerStatisticsNaturalCommunication2018}. Alternatively, they
could serve as a signal of stress that is simply caused by extended chasing
without influencing the behavior of the winner. If chirps communicated
submission, leading to offset of chasing, the prediction would be that chasing
events were chirps were produced should have been shorter. This is not the
case, as the opposite could be observed: Chasing events, in which losers
produced chirps lasted significantly longer (p<0.0001, d=19.3). In every trial
in which the loser produced chirps during chasing, chirps were produced when
the median of the chasing durations was increased (figure \ref{fig:stats} B). 
An influence of chirping on the latency between the current and the next
chasing could not be observed. This suggests that competition losers started to
produce chirps as a response to beeing continuously chased by the competition
winner.

To understand with which particular behavior the chirps produced by the
competition losers or winners correlated temporally, I approximated how the
chirp rate changes during these events. As already implied in figure
\ref{fig:competition_trial_overview}, chirping of the competition loser but not
the competition winner was closely tied to chasing events (figure
\ref{fig:stats} C \& D vs E \& F). Chirp rates of the competition loser
significantly increased at the onset of being chased by the competition winner
(figure \ref{fig:stats} E). This effect is even more pronounced during the time
at which the chasing seizes: Here, the chirp rate of the competition loser
almost tripled and remained elevated for approximately \SI{10}{\second} beyond
the offset (figure \ref{fig:stats} F). This effect could not be observed for
competition winners, which produced little chirps and if so, showed no
correlations with the onsets or offsets of chasing events (figure
\ref{fig:stats} C, D).

\begin{figure}
  \begin{center}
    \includegraphics[width=\textwidth]{../figures/competition_trial_overview.pdf}
  \end{center}
  \caption{
    Overview of a single competition trial. The spectrogram illustrates the
    evolution of the discharge frequencies of both fish. Fish were left to
    compete for three hours while the light was off. This is also the time in
    which most of the frequency modulations are visible on the spectrogram. As
    the light is turned off, frequency changes of either fish become much
    rarer. Event plots show the times of the two annotated spatio-temporal
    interactions between the fish, the chasing and contact events. The majority
    of agonistic events are chasing events. Below are the event plots for the
    communication signals produced by either fish, separated by the competition
    status of the respective fish. Most rises but also most chirps are produced
    by the fish that lost the competition and did not get access to the
    shelter. Additionally, bouts of subsequent chirps produced by the loser
    visibly correlate with the chasing events.
  }
  \label{fig:competition_trial_overview}
  \vspace{0.5cm}
\end{figure}

\begin{figure}
  \begin{center}
    \includegraphics[width=\textwidth]{../figures/plot_behavioral_correlations_of_chirps.pdf}
  \end{center}
  \caption{
    Statistical summary of the role of loser chirps during ritualized chasing
    events. Points are connected according to the unique fish pairing in which
    they were observed. Competition losers produced more chirps than the
    competition winners (A). If the loser produced a chirp, chasing events
    lasted longer (B). Most chirps of the loser are produced close to
    the offset of the chasing event (figure \ref{fig:stats} F), which suggests
    that it is more likely that being chased for a longer time causes the
    losers to chirp. Chirp rates of competition losers centered around the
    onset and the offset of chasing events respectively (C-F). The chirp rate
    was estimated by convolving a Gaussian kernel with a binary chirp train.
    The confidence intervals are estimated by jackknifing and the baseline by
    bootstrap resampling. Chirp rates of the competition losers significantly
    increase after a chasing events starts (E). Specifically, most of these
    chirps that increase the chirp rate are precisely produced around the
    offset of chasing events (F), where the chirp rate almost triples.
  }
  \label{fig:stats}
  \vspace{0.5cm}
\end{figure}

% chapter Results (end)
